\documentclass[11pt]{article}
\usepackage[T1]{fontenc}
\usepackage[utf8]{inputenc}
\usepackage{lmodern}
\usepackage[a4paper,margin=1in]{geometry}
\usepackage{microtype}
\usepackage{amsmath,amssymb,amsthm,mathtools}
\usepackage{booktabs,array}
\usepackage[dvipsnames]{xcolor}
\usepackage[colorlinks=true,linkcolor=MidnightBlue,citecolor=MidnightBlue,
            urlcolor=MidnightBlue]{hyperref}
\usepackage{enumitem}
\usepackage{setspace}
\onehalfspacing
\setlist[itemize]{topsep=4pt,itemsep=2pt,leftmargin=1.5em}
\setlist[enumerate]{topsep=4pt,itemsep=2pt,leftmargin=1.8em}

\newtheorem{theorem}{Theorem}[section]
\newtheorem{lemma}[theorem]{Lemma}
\newtheorem{proposition}[theorem]{Proposition}
\newtheorem{corollary}[theorem]{Corollary}
\theoremstyle{definition}
\newtheorem{definition}[theorem]{Definition}
\newtheorem{construction}[theorem]{Construction}
\theoremstyle{remark}
\newtheorem{remark}[theorem]{Remark}

\newcommand{\DR}{\mathsf{DR}}
\newcommand{\PO}{\mathsf{PO}}
\newcommand{\PP}{\mathsf{PP}}
\newcommand{\PPI}{\mathsf{PPI}}
\newcommand{\EQ}{\mathsf{EQ}}
\newcommand{\Base}{\{\DR,\PO,\EQ,\PP,\PPI\}}
\newcommand{\BaseEdge}{\{\DR,\PO,\PP,\PPI\}}
\newcommand{\cl}{\operatorname{cl}}
\newcommand{\Tp}{\operatorname{Tp}}
\newcommand{\tp}{\operatorname{tp}}
\newcommand{\md}{\operatorname{md}}
\newcommand{\st}{\operatorname{st}}
\newcommand{\comp}{\operatorname{comp}}
\newcommand{\inv}{\operatorname{inv}}
\newcommand{\orig}{\operatorname{orig}}
\newcommand{\Need}{\operatorname{Need}}
\newcommand{\Safe}{\operatorname{Safe}}
\newcommand{\Core}{\mathsf{Core}}
\newcommand{\Out}{\mathsf{Out}}
\newcommand{\Front}{\mathsf{Front}}
\newcommand{\Wit}{\operatorname{Wit}}
\newcommand{\Iface}{\operatorname{Iface}}
\newcommand{\Child}{\operatorname{Child}}
\newcommand{\Par}{\operatorname{Par}}
\newcommand{\Dom}{\mathsf{Dom}}
\newcommand{\Mos}{\mathsf{Mos}}
\newcommand{\gen}{\mathsf{gen}}
\newcommand{\eqcong}{\equiv_{\EQ}}
\newcommand{\reach}{\rightsquigarrow}
\newcommand{\ALCIRCC}[1]{\mathcal{ALCI}_{\mathsf{RCC#1}}}

\title{\textbf{Completeness Extraction for $\ALCIRCC{5}$ (v2, Repaired)}\\[0.4em]
\large Closing the model-to-quotient gap via witness-generated submodels,\\
generated covers, and reachability discharge}
\author{Claude (Anthropic, Opus 4.7)\\
\small Prompted by Michael Wessel}
\date{May 2026}

\begin{document}
\maketitle

\begin{center}
\fbox{\parbox{0.92\textwidth}{%
\textbf{Disclaimer.}\;
This document was generated by Claude (Anthropic), prompted by Michael
Wessel.  The results have not been independently verified.  We invite
scrutiny and corrections.}}
\end{center}

\begin{abstract}
This is a repaired version of the completeness-extraction note for
the split-forest approach to $\ALCIRCC{5}$ decidability.  An
earlier version~\cite{CompletenessV1} contained two structural
defects that prevented it from establishing the model-to-quotient
direction: a Hasse construction that assumed a well-founded ambient
$\le_{\PP}$ DAG with a $\le_{\PP}$-minimal root (which fails for
concepts forcing infinite $\PP$-chains, e.g.\ $\exists\PP.\top
\sqcap\forall\PP.\exists\PP.\top$), and a triple-coherence proof
with witness conflation across distinct occurrences.

The repair follows the structural ideas of GPT~5.5's coherent
generated split-forest proof~\cite{GPT55Generated}, but \emph{without}
introducing B\"uchi automata or $\omega$-acceptance machinery.  We
apply seven concrete repairs:

\begin{enumerate}[label=(R\arabic*),leftmargin=2em]
\item Pass to a \emph{witness-generated submodel} (sparse,
  logic-generated, no density issue);
\item Use \emph{generated cover edges} (selected witness-cover edges
  in the generated submodel), not semantic Hasse covers;
\item Drop the ambient-root requirement; permit \emph{self-loops on
  the parent function} of the finite quotient, so ultimately
  periodic ancestor trunks are first-class;
\item Make the \emph{containment orientation} explicit (parent =
  proper superpart, child = proper part);
\item Replace raw pair-tables with \emph{support-closed mosaics} that
  store joint atomic networks, fixing the C5 witness conflation;
\item Add an explicit \emph{typed-$\EQ$ congruence} axiom, separate
  from blocking equivalence;
\item Discharge transitive $\PP/\PPI$ eventualities by
  \emph{finite-graph reachability} on the finite quotient.
\end{enumerate}

The crucial observation is that the cycle-realization condition
of~\cite{GPT55Generated} reduces to ordinary reachability over the
finite quotient: $\exists\PP.D \in \lambda(\sigma)$ requires
$D \in \lambda(\sigma')$ for some $\sigma'$ reachable from $\sigma$
along parent edges (including self-loops).  This is decidable in
$O(|S|^2)$ time per eventuality and requires no $\omega$-language
machinery.

Combined with GPT-5.4's soundness chain~\cite{SplitForest}, this
yields decidability of concept satisfiability in $\ALCIRCC{5}$
under strong $\EQ$ semantics.
\end{abstract}


% ═══════════════════════════════════════════════════════════════
\section{Introduction}
\label{sec:intro}
% ═══════════════════════════════════════════════════════════════

The split-forest proposal for $\ALCIRCC{5}$
decidability~\cite{SplitForest} proves the soundness direction in
full detail: every valid finite rank-$d$ quotient unfolds into a
genuine strong-$\EQ$ model via canonical refinements of disjunctive
RCC5 networks.  The completeness direction---that every satisfiable
concept admits a valid finite quotient---was stated as a proof
sketch.

A first attempt to write out this completeness extraction
(``v1'')~\cite{CompletenessV1} appeared in April 2026.  Subsequent
review by GPT-5.5 identified two structural defects:

\paragraph{Defect 1: Hasse construction with rooted-tree assumption.}
v1 began by orienting the proper-part relation $\le_{\PP}$ as a DAG
and choosing a $\le_{\PP}$-minimal element as ambient root.  This
works when the model has a unique minimum or when one can be added
freshly.  But the concept
\[
C\!\uparrow \;:=\; \exists\PP.\top \;\sqcap\; \forall\PP.\exists\PP.\top
\]
forces a strictly increasing $\PP$-chain $d_0\PP d_1\PP d_2\PP\cdots$
with no $\le_{\PP}$-minimum.  Adding a fresh top via
$\rho(d_\top,d)=\PP$ \emph{inverts} the orientation: by RCC5
semantics, $d_\top\PP d$ means $d_\top$ is a proper part of $d$, not
the other way around.  v1's construction therefore breaks on any
concept forcing an infinite $\PP$-chain.

\paragraph{Defect 2: Witness conflation in (C5) triple coherence.}
v1's proof of (C5) implicitly combined witnesses from different
occurrences as if they formed a coherent triple.  Raw pair tables
admit independent witnesses for $M(\chi_1,\chi_2)$ and
$M(\chi_1,\chi_3)$ that may not be jointly realizable.  Triple
coherence requires a single triple $(x_1,x_2,x_3)$ realizing all
three relations simultaneously; raw tables do not enforce this.

\medskip
GPT-5.5's coherent generated split-forest
proof~\cite{GPT55Generated} repairs both defects.  We follow its
approach with one important difference: where GPT-5.5 phrases the
eventuality discharge in automata-theoretic language
(``cycle-realization condition'', Definition~6.4 of that paper),
we observe that this condition reduces to \emph{finite-graph
reachability over the finite quotient}, and we present the proof
in that finitary style.

\subsection{Plan}

Section~\ref{sec:prelim} recaps types, closure, need families, and
the model-edge relation-safety lemma.
Section~\ref{sec:witness-generated} introduces the
witness-generated submodel reduction.
Section~\ref{sec:generated-covers} fixes the containment
orientation and defines generated cover edges.
Section~\ref{sec:split-tree} constructs the generated split-tree
presentation, with no ambient root and self-loops permitted on the
parent function.
Section~\ref{sec:rank-d} defines rank-$d$ profiles and the finite
quotient.
Section~\ref{sec:mosaics} extracts support-closed sibling-interface
mosaics.
Section~\ref{sec:witness-menu} defines witness menus for $\DR/\PO$
existentials.
Section~\ref{sec:typed-eq} states the typed-$\EQ$ congruence axiom.
Section~\ref{sec:coherence} verifies (C1)--(C5).
Section~\ref{sec:reachability} proves reachability discharge for
$\PP/\PPI$ eventualities.
Section~\ref{sec:validity} assembles validity.
Section~\ref{sec:main} states the main theorem and decidability.
Section~\ref{sec:discussion} compares with GPT-5.5's automata
version and discusses honest scope.

\subsection{Main result}

\begin{theorem}[Completeness extraction]\label{thm:main}
If $\mathcal{I},d_\ast\models C_0$ in $\ALCIRCC{5}$ under strong
$\EQ$ semantics, then there exists a valid finite generated
split-forest quotient $Q$ for $C_0$ with $d=\md(C_0)$.
\end{theorem}

Combined with the soundness chain of~\cite{SplitForest,GPT55Generated},
this yields:

\begin{corollary}[Decidability of $\ALCIRCC{5}$]\label{cor:decidability}
Concept satisfiability in $\ALCIRCC{5}$ under strong $\EQ$ semantics
is decidable.
\end{corollary}


% ═══════════════════════════════════════════════════════════════
\section{Preliminaries}
\label{sec:prelim}
% ═══════════════════════════════════════════════════════════════

We work under the conventions of~\cite{SplitForest}.  This section
recalls only the notions used in the extraction.

\subsection{Types and closure}

Fix a concept $C_0$ in negation normal form (NNF).  The
Fischer-Ladner closure $\cl(C_0)$ is the smallest set containing
$C_0$ closed under subformulas and NNF complements.  A
\emph{Hintikka type} $\tau \in \Tp(C_0) \subseteq 2^{\cl(C_0)}$ is
a maximally consistent subset satisfying the standard propositional
closure conditions.  For $a$ in a model $\mathcal{I}$, write
$\tp_{\mathcal{I}}(a) = \{D \in \cl(C_0) : \mathcal{I},a\models D\}$;
this is a Hintikka type.

\subsection{Need families and relation-safety}

\begin{definition}[Per-relation needs]\label{def:need}
For $\tau \in \Tp(C_0)$ and $R \in \BaseEdge$, define
\[
\Need_R(\tau) := \{D \in \cl(C_0) : \forall R.D \in \tau\}.
\]
A typed edge $(\tau, R, \sigma)$ with $\tau,\sigma\in\Tp(C_0)$ is
\emph{relation-safe} when
\[
\Need_R(\tau) \subseteq \sigma
\quad\text{and}\quad
\Need_{\inv(R)}(\sigma) \subseteq \tau.
\]
\end{definition}

\begin{lemma}[Model edges are relation-safe]\label{lem:model-safe}
In any model $\mathcal{I}$ of $\ALCIRCC{5}$, if $\rho(a,b) = R \in
\BaseEdge$ for distinct elements $a,b$, then $(\tp_{\mathcal{I}}(a),
R, \tp_{\mathcal{I}}(b))$ is relation-safe.
\end{lemma}

\begin{proof}
If $\forall R.D \in \tp_{\mathcal{I}}(a)$, then $\mathcal{I},a
\models \forall R.D$ and $\rho(a,b)=R$, hence $\mathcal{I},b
\models D$, so $D \in \tp_{\mathcal{I}}(b)$.  The dual inclusion
follows by applying the same argument with the inverse relation
at $b$.
\end{proof}

This elementary observation is the engine of the entire extraction:
every relation realized in the model is automatically relation-safe
for its endpoints' types.

\subsection{RCC5 composition}

We use the standard RCC5 composition table $\comp(R,S) \subseteq
\Base$ (see~\cite[Table~1]{GPT55Generated}).  The composition is
sound for any abstract RCC5 frame: $\rho(a,c) \in
\comp(\rho(a,b),\rho(b,c))$.


% ═══════════════════════════════════════════════════════════════
\section{Witness-generated submodel reduction}
\label{sec:witness-generated}
% ═══════════════════════════════════════════════════════════════

The first repair is to restrict attention to sparse,
logic-generated submodels.  This sidesteps every concern about
arbitrary-density semantic orders and lets us define cover edges
relative to a finite generated set.

\begin{definition}[Selected witness policy]\label{def:witness-policy}
Let $\mathcal{I}$ be a model.  A \emph{selected witness policy} $w$
chooses, for every $a \in \Delta^{\mathcal{I}}$ and every
existential $\exists R.D_0 \in \cl(C_0)$ with $a\in (\exists
R.D_0)^{\mathcal{I}}$, one element
\[
w(a, \exists R.D_0) \in \Delta^{\mathcal{I}}
\quad\text{with}\quad
\rho(a, w(a,\exists R.D_0)) = R \text{ and } w(a,\exists R.D_0) \in
D_0^{\mathcal{I}}.
\]
We write $a \xrightarrow{R}_w b$ when $b$ is so selected.
\end{definition}

\begin{definition}[Witness-generated submodel]
\label{def:witness-generated}
Let $a_0 \in \Delta^{\mathcal{I}}$.  The \emph{witness-generated
domain} $D_w(a_0) \subseteq \Delta^{\mathcal{I}}$ is the smallest
set containing $a_0$ and closed under all selected witnesses for
closure-existentials true at its members.  The \emph{induced
submodel} $\mathcal{I}_w$ has domain $D_w(a_0)$ with
$\rho^{\mathcal{I}_w}$ and atomic concept interpretations restricted
from $\mathcal{I}$.
\end{definition}

\begin{lemma}[Witness-generated submodel
lemma]\label{lem:witness-generated}
If $a_0 \in C_0^{\mathcal{I}}$, then for every $a \in D_w(a_0)$ and
every $E \in \cl(C_0)$,
\[
a \in E^{\mathcal{I}} \;\Longleftrightarrow\; a \in E^{\mathcal{I}_w}.
\]
In particular, $a_0 \in C_0^{\mathcal{I}_w}$.
\end{lemma}

\begin{proof}
The induced substructure remains a strong abstract RCC5 frame: all
RCC5 frame requirements (irreflexivity of $\BaseEdge$, converse,
composition) are universal over retained pairs and triples and are
preserved under restriction.  Atomic concepts are preserved by
restriction.

Boolean cases of the equivalence are immediate from the inductive
hypothesis.

For $E = \exists R.F$: if $a \in E^{\mathcal{I}_w}$, retention into
$\mathcal{I}_w$ implies the witness exists in $\mathcal{I}$, so
$a \in E^{\mathcal{I}}$.  Conversely, if $a \in E^{\mathcal{I}}$,
the selected witness $b = w(a,\exists R.F)$ lies in $D_w(a_0)$ and
satisfies $F$ in $\mathcal{I}$; by induction $b \in
F^{\mathcal{I}_w}$, so $a \in E^{\mathcal{I}_w}$.

For $E = \forall R.F$: if $a \in E^{\mathcal{I}}$, every retained
$R$-neighbour of $a$ is also an $R$-neighbour in $\mathcal{I}$ and
satisfies $F$; by induction it satisfies $F$ in $\mathcal{I}_w$.
Conversely, suppose $a \notin E^{\mathcal{I}}$.  Then $a \in
(\exists R.\overline{F})^{\mathcal{I}}$, where $\overline{F}$ is
the NNF complement, and $\exists R.\overline{F} \in \cl(C_0)$
because the closure is closed under NNF complements.  The selected
witness for $\exists R.\overline{F}$ is retained and, by induction,
satisfies $\overline{F}$ in $\mathcal{I}_w$; hence $a \notin
E^{\mathcal{I}_w}$.
\end{proof}

\begin{corollary}[Sparse generated models suffice]\label{cor:sparse}
Every satisfiable $C_0$ has a witness-generated model with $|D_w(a_0)|
\le \aleph_0$.  The cover-tree extraction below need only address
this sparse class.
\end{corollary}

\begin{remark}[Why the density issue dissolves]
Lemma~\ref{lem:witness-generated} replaces v1's flawed Hasse
construction with a model-theoretic thinning.  The resulting
$\mathcal{I}_w$ contains exactly the elements that the logic
witnesses for: the evaluation point and selected witnesses for
existentials in the closure.  Density of an arbitrary semantic
$\PP/\PPI$ order is irrelevant, because we never quantify over it.
\end{remark}


% ═══════════════════════════════════════════════════════════════
\section{Generated cover edges and containment orientation}
\label{sec:generated-covers}
% ═══════════════════════════════════════════════════════════════

We now fix the orientation convention used throughout the rest of
the paper.

\begin{quote}
\textbf{Containment orientation.}\;
\emph{Parent} = proper superpart; \emph{child} = proper part.\\
Equivalently: along a tree branch, $\text{child}\,\PP\,\text{parent}$
and $\text{parent}\,\PPI\,\text{child}$.
\end{quote}

This convention determines which existentials produce ``upward''
moves and which produce ``downward'' moves in the split tree.
$\exists\PP.D$ at $a$ requires a proper-superpart witness, hence an
\emph{ancestor} in the tree.  $\exists\PPI.D$ at $a$ requires a
proper-part witness, hence a \emph{descendant}.

\begin{definition}[Generated cover edge]\label{def:generated-cover}
Let $A \subseteq D_w(a_0)$ be a finite witness-generated subset of
the submodel.  A \emph{generated $\PPI$-cover edge}
$x \xrightarrow{\PPI}_{\gen,A} y$ holds when:
\begin{itemize}
\item $x, y \in A$ and $\rho(x,y) = \PPI$ (i.e., $y$ is a proper
  part of $x$);
\item no element $z \in A$ satisfies $\rho(x,z) = \PPI$ and
  $\rho(z,y) = \PPI$ (no selected element of $A$ lies strictly
  between $x$ and $y$ in the proper-part order).
\end{itemize}
The dual generated $\PP$-cover edge $x \xrightarrow{\PP}_{\gen,A} y$
is defined symmetrically.
\end{definition}

\begin{remark}[Generated covers are not Hasse covers]
A generated cover edge is immediate \emph{relative to the finite
selected set $A$}.  It is not claimed to be an immediate cover in
the full ambient $\PP$-order.  Repaired wording from the survival
note~\cite{GPT55Survival}:
\begin{quote}
\emph{A cover edge is a generated witness-cover edge, not a
semantic Hasse cover.}
\end{quote}
Constructing covers relative to $A$ is well-defined even when the
ambient order is dense or has no Hasse covers at all; we simply
pick the immediate-in-$A$ relation.  Crucially, $A$ grows with the
generation depth, but at every finite depth the relative cover
relation is well-defined.
\end{remark}


% ═══════════════════════════════════════════════════════════════
\section{Generated split-tree presentation}
\label{sec:split-tree}
% ═══════════════════════════════════════════════════════════════

We construct a (possibly infinite) split tree from a witness-generated
model.  No ambient root is required.

\begin{construction}[Generated split-tree]\label{con:split-tree}
Let $\mathcal{I}_w$ be a witness-generated model with evaluation
point $a_0$.  Build a tree $T$ of \emph{occurrences} together with
maps
\[
\orig : T \to D_w(a_0), \quad
\lambda : T \to \Tp(C_0), \quad
\mathsf{par} : T \rightharpoonup T,
\]
inductively as follows.

\begin{enumerate}
\item \textbf{Seed.}\; Create a root occurrence $u_0$ with
$\orig(u_0) = a_0$ and $\lambda(u_0) = \tp_{\mathcal{I}_w}(a_0)$.

\item \textbf{Upward expansion (PP-witnesses).}\; For each
occurrence $u$ and each $\exists\PP.D \in \lambda(u)$:
\begin{itemize}
\item If $\mathsf{par}(u)$ is undefined or
  $\lambda(\mathsf{par}(u)) \not\ni D$, select a witness $b =
  w(\orig(u),\exists\PP.D) \in D_w(a_0)$ with $\rho(\orig(u),b) =
  \PP$ and create a new occurrence $u^\uparrow$ with
  $\orig(u^\uparrow) = b$, $\lambda(u^\uparrow) =
  \tp_{\mathcal{I}_w}(b)$, $\mathsf{par}(u) = u^\uparrow$.
\item If a parent already records $D$, no new occurrence is needed.
\end{itemize}

\item \textbf{Downward expansion (PPI-witnesses).}\; For each
occurrence $u$ and each $\exists\PPI.D \in \lambda(u)$, select a
witness $b = w(\orig(u),\exists\PPI.D)$ and create a child
occurrence $u^\downarrow$ with $\orig(u^\downarrow) = b$,
$\lambda(u^\downarrow) = \tp_{\mathcal{I}_w}(b)$, and
$\mathsf{par}(u^\downarrow) = u$.

\item \textbf{Sibling expansion (DR/PO-witnesses).}\; For each
occurrence $u$ and each $\exists R.D \in \lambda(u)$ with $R \in
\{\DR,\PO\}$, select a witness $b = w(\orig(u),\exists R.D)$ and
record an outgoing $R$-witness edge from $u$ to a node mapped to
$b$ (typically a sibling of $u$ under a common ancestor).

\item \textbf{Splitting.}\; If a single element $b \in D_w(a_0)$
must appear under multiple distinct parents in this construction,
introduce one occurrence per containment branch.  These
occurrences are declared $\EQ$-mates (typed equivalence classes
under the relation defined in Section~\ref{sec:typed-eq}).

\item \textbf{Edge labels.}\; For comparable occurrences $u, v$ on
the same branch with $v = \mathsf{par}^k(u)$ for some $k > 0$, set
$\rho_T(u,v) = \PP$ and $\rho_T(v,u) = \PPI$ (transitively closed
along the branch).  For incomparable occurrences, $\rho_T(u,v)$ is
determined by the realized RCC5 relation $\rho(\orig(u),\orig(v))$
in $\mathcal{I}_w$.
\end{enumerate}
\end{construction}

\begin{remark}[Branches need not terminate]\label{rem:no-root}
The parent chain of an occurrence may be:
\begin{itemize}
\item finite, terminating at an occurrence with no $\exists\PP.D$
  in its type;
\item one-sided infinite (e.g., for $C\!\uparrow$);
\item bi-infinite (e.g., for concepts forcing both unbounded $\PP$
  and unbounded $\PPI$ chains).
\end{itemize}
We do \emph{not} require a $\le_{\PP}$-minimal element or any
ambient top.  In the finite quotient (Section~\ref{sec:rank-d}),
infinite parent chains are represented by self-loops or longer
cycles on the parent function.
\end{remark}

\begin{lemma}[Existence]\label{lem:split-extraction}
Every witness-generated model $(\mathcal{I}_w, a_0)$ admits a
generated split tree as in Construction~\ref{con:split-tree}.
\end{lemma}

\begin{proof}
Construct occurrences coinductively.  At every step, the existing
occurrences are finite (since each step adds finitely many
witnesses for the closure-existentials of one type).  The
construction is well-defined because $\cl(C_0)$ is finite, and
there are only finitely many distinct closure-existentials at any
type.  Pull-back of $\rho$ via $\orig$ inherits converse and
composition.  Type assignments inherit Hintikka conditions from
$\mathcal{I}_w$ via Lemma~\ref{lem:witness-generated}.
\end{proof}


% ═══════════════════════════════════════════════════════════════
\section{Rank-$d$ profiles and the finite quotient}
\label{sec:rank-d}
% ═══════════════════════════════════════════════════════════════

We now define the finite combinatorial object that summarizes the
(possibly infinite) split tree.

\begin{definition}[Profile components]\label{def:profile-components}
For an occurrence $u \in T$ at rank $k$, the \emph{rank-$k$
profile} $\st_k(u)$ records:
\begin{enumerate}
\item the closure type $\lambda(u)$;
\item the \emph{typed-$\EQ$ colour} $c_{\EQ}(u)$ (a label from a
  finite alphabet, used to track $\EQ$-mate classes);
\item the rank-$(k-1)$ profile of $\mathsf{par}(u)$ (with $\bot$ if
  undefined);
\item the rank-$(k-1)$ \emph{child multiplicities}: a map
  $\Sigma_{k-1} \to \{0,1,2{+}\}$ counting children by their rank-$(k-1)$
  profile;
\item the rank-$(k-1)$ \emph{incident sibling-interface mosaics}
  for ordered pairs of children
  (Section~\ref{sec:mosaics});
\item the rank-$(k-1)$ \emph{witness menu} for $\DR/\PO$
  existentials (Section~\ref{sec:witness-menu});
\item the \emph{vertical eventuality summary}: for each
  $\exists\PP.D \in \lambda(u)$, the set of ranks at which $D$
  appears in the parent chain; symmetric for $\exists\PPI.D$.
\end{enumerate}
A \emph{rank-$0$ profile} is just a Hintikka type.  The finite
profile alphabet $\Sigma_d$ at rank $d = \md(C_0)$ has bounded
cardinality (Lemma~\ref{lem:finite-index}).
\end{definition}

\begin{lemma}[Finite-index]\label{lem:finite-index}
For each $k \le d$, the set $\Sigma_k$ of rank-$k$ profiles is
finite, and $|\Sigma_k|$ is bounded by an effectively computable
function of $|\cl(C_0)|$.
\end{lemma}

\begin{proof}
Identical to the finite-index argument
in~\cite[Lemma~7.4]{GPT55Generated}: closure types are finite;
typed-$\EQ$ colours are drawn from a finite palette of size
$O(2^{|\cl(C_0)|})$; child multiplicities are truncated at $2{+}$
(threshold counting); incident mosaics are bounded sets of bounded
networks over a finite endpoint-summary alphabet
(Section~\ref{sec:mosaics}); witness-menu entries are pairs from a
finite set; eventuality summaries are subsets of $\cl(C_0)$.
Induction on $k$ closes the bound.
\end{proof}

\begin{definition}[Finite quotient]\label{def:quotient}
The \emph{finite generated split-forest quotient} of $T$ is
\[
Q = (S, \lambda_Q, \mathsf{par}_Q, \Child_Q, \Iface_Q, \Wit_Q,
c_{\EQ,Q}, \sigma_\ast)
\]
where:
\begin{itemize}
\item $S = \{\st_d(u) : u \in T\} \subseteq \Sigma_d$;
\item $\lambda_Q(\sigma)$ is the closure type at $\sigma$;
\item $\mathsf{par}_Q : S \rightharpoonup S$ is the projected parent
  function: $\mathsf{par}_Q(\sigma) = \st_d(\mathsf{par}(u))$ for
  any $u \in T$ with $\st_d(u) = \sigma$ (well-defined by profile
  equality, including the parent slot; $\mathsf{par}_Q(\sigma) =
  \sigma$ is permitted, as is $\mathsf{par}_Q(\sigma) = \bot$);
\item $\Child_Q, \Iface_Q, \Wit_Q$ are inherited from the profile
  components;
\item $c_{\EQ,Q}$ is the projected typed-$\EQ$ colouring;
\item $\sigma_\ast = \st_d(u_\ast)$ for the occurrence $u_\ast$
  with $\orig(u_\ast) = a_0$.
\end{itemize}
\end{definition}

\begin{remark}[Self-loops are allowed]\label{rem:self-loops}
The parent function $\mathsf{par}_Q$ may have self-loops:
$\mathsf{par}_Q(\sigma) = \sigma$.  This represents an infinite
upward $\PP$-chain that, at rank $d$, has stabilized to a single
profile.  The unfolding interprets each traversal of the self-loop
as a fresh occurrence.  Concretely, for $C\!\uparrow = \exists\PP.\top
\sqcap\forall\PP.\exists\PP.\top$, the quotient has a single
profile $\sigma$ with $\mathsf{par}_Q(\sigma) = \sigma$ and
$\lambda_Q(\sigma) \ni \exists\PP.\top$, $\lambda_Q(\sigma) \ni
\forall\PP.\exists\PP.\top$.

\textbf{Crucially, $\mathsf{par}_Q(\sigma) = \sigma$ does \emph{not}
mean ``$\sigma$ is its own parent in the model.''}  It means: the
parent occurrence has the same rank-$d$ profile.  In the unfolding,
every traversal produces a distinct semantic occurrence.  This is
a blocking cycle, not an $\EQ$-identification.
\end{remark}

\begin{lemma}[Finiteness of Q]\label{lem:finite-Q}
$|S| \le |\Sigma_d| < \infty$.
\end{lemma}


% ═══════════════════════════════════════════════════════════════
\section{Sibling-interface descriptors as support-closed mosaics}
\label{sec:mosaics}
% ═══════════════════════════════════════════════════════════════

This section replaces v1's raw pair-tables with support-closed
mosaics, fixing the C5 witness conflation.

\subsection{Status partition}

For an occurrence $u$ with two distinct child subtrees $T_s, T_t$
at children $s, t$, we partition descendants of $s$ by their RCC5
relation to $t$:

\begin{definition}[Status partition]\label{def:status}
For $x$ a descendant of $s$ (including $s$ itself):
\begin{itemize}
\item $\sst{}(x) = \Core$ if $\orig(x)$ lies below both $\orig(s)$
  and $\orig(t)$ in the generated $\PPI$-order, or is an
  $\EQ$-mate visible to both sides;
\item $\sst{}(x) = \Out$ if $x \notin \Core$ and
  $\rho(\orig(x),\orig(t)) = \DR$;
\item $\sst{}(x) = \Front$ if $x \notin \Core$ and
  $\rho(\orig(x),\orig(t)) = \PO$.
\end{itemize}
\end{definition}

\begin{lemma}[Three-way partition]\label{lem:status-partition}
Every descendant $x$ of $s$ has exactly one status in
$\{\Core, \Out, \Front\}$.  Symmetrically for descendants of $t$.
\end{lemma}

\begin{proof}
By Lemma~\ref{lem:open-cross} and the open cross-edge restriction
of~\cite[Theorem~1.4]{SplitForest}: if $x$ is a non-core descendant
of $s$, then $\rho(\orig(x),\orig(t)) \in \{\DR,\PO\}$ (it cannot
be $\PP$, $\PPI$, or $\EQ$ without forcing $x$ into $\Core$ via
generated containment closure).
\end{proof}

\begin{lemma}[Open cross-edge]\label{lem:open-cross}
If $x$ is a non-core descendant of $s$ and $y$ is a non-core
descendant of $t$, then $\rho(\orig(x),\orig(y)) \in \{\DR,\PO\}$.
\end{lemma}

\begin{proof}
$\EQ$ would force both into the typed-$\EQ$ congruence and hence
into $\Core$.  $\PP$ would force $x$ into the generated
$\PP$-closure of material below $\orig(t)$, hence visible to $t$,
contradicting non-coreness.  $\PPI$ is symmetric.
\end{proof}

\subsection{Endpoint summaries and atomic mosaics}

\begin{definition}[Endpoint summary]\label{def:endpoint}
An \emph{endpoint summary} at rank $k$ consists of:
\begin{enumerate}
\item a rank-$k$ profile $p \in \Sigma_k$;
\item a status $\sigma \in \{\Core,\Out,\Front\}$;
\item a typed-$\EQ$ colour $c_{\EQ} \in \mathsf{Colours}$.
\end{enumerate}
The set $E_k$ of endpoint summaries is finite.
\end{definition}

\begin{definition}[Atomic support mosaic]\label{def:mosaic}
An \emph{atomic support mosaic} on a finite set $V$ of vertices,
each labelled by an endpoint summary, is a complete labelling
$\mu : V \times V \to \Base$ such that:
\begin{enumerate}[label=(M\arabic*),leftmargin=2em]
\item $\mu(v,v) = \EQ$ for all $v$;
\item $\mu(w,v) = \inv(\mu(v,w))$;
\item $\mu(u,w) \in \comp(\mu(u,v),\mu(v,w))$ (RCC5 composition);
\item every edge is relation-safe for the closure types of its
  endpoint summaries (Definition~\ref{def:need});
\item $\Out$-status endpoints have only $\DR$ as cross-status
  outgoing label;
\item non-core $\Front$--$\Front$ cross-status edges have labels in
  $\{\DR,\PO\}$;
\item endpoints with the same typed-$\EQ$ colour have label $\EQ$
  (typed-$\EQ$ congruence respects the mosaic).
\end{enumerate}
\end{definition}

\begin{definition}[Sibling descriptor]\label{def:descriptor}
A \emph{rank-$k$ sibling descriptor} for an ordered child pair
$(s,t)$ is a finite set $\Mos_k(s,t)$ of atomic support mosaics
over endpoint summaries at rank $k$, capturing the joint atomic
networks that occur in the model among finitely many endpoints
in $T_s \cup T_t \cup \{s, t\}$.

The descriptor is \emph{support-closed} when it is closed under:
\begin{enumerate}[label=(SC\arabic*),leftmargin=2em]
\item \textbf{Restriction.}\; Every induced submosaic on a
  vertex subset is in $\Mos_k(s,t)$.
\item \textbf{Triple composition.}\; If a 3-vertex submosaic on
  $(v_1,v_2,v_3)$ has labels $(R_{12}, R_{13}, R_{23})$, then
  $R_{13} \in \comp(R_{12},R_{23})$.
\item \textbf{Typed-$\EQ$ replacement.}\; Replacing a vertex by an
  $\EQ$-mate of the same endpoint summary yields a mosaic in
  $\Mos_k(s,t)$.
\item \textbf{One-step containment extension.}\; Adding a generated
  $\PPI$-child of a vertex (with appropriate label propagation
  by composition with $\PP$ or $\PPI$) yields a mosaic in
  $\Mos_k(s,t)$.
\end{enumerate}
\end{definition}

\begin{lemma}[Support closure is finite]\label{lem:support-closure}
For fixed rank $k$, the support closure of any finite set of
mosaics is finite and effectively computable.
\end{lemma}

\begin{proof}
The endpoint summary alphabet $E_k$ is finite.  Mosaics are bounded
in size by some constant depending only on $|\cl(C_0)|$ (we may
restrict to mosaics on at most $|\cl(C_0)| + 4$ vertices, since
larger mosaics can be decomposed into overlapping smaller ones via
restriction).  Each closure operation adds elements from this
finite universe; iteration terminates.
\end{proof}

\subsection{Extraction from a model}

\begin{construction}[Mosaic extraction]\label{con:mosaic-extract}
For each ordered child pair $(s,t)$ in the split tree at rank $k$:

\begin{enumerate}
\item For every finite vertex set $V \subseteq T_s \cup T_t \cup
  \{s,t\}$, form the \emph{realized mosaic} $\mu_V$ with $\mu_V(v,w)
  = \rho(\orig(v),\orig(w))$ and endpoint summaries given by the
  rank-$k$ profile, status, and typed-$\EQ$ colour of each $v \in
  V$.
\item Initialize $\Mos_k^{\mathrm{raw}}(s,t)$ as the set of
  realized mosaics for $V$ ranging over finite subsets up to a
  bounded vertex count.
\item Close under (SC1)--(SC4) of Definition~\ref{def:descriptor}.
\end{enumerate}
The result $\Mos_k(s,t)$ is the \emph{extracted sibling descriptor}.
\end{construction}

\begin{lemma}[Extracted descriptors are coherent]
\label{lem:descriptor-coherent}
The descriptor $\Mos_k(s,t)$ extracted in
Construction~\ref{con:mosaic-extract} satisfies (M1)--(M7) and
(SC1)--(SC4).
\end{lemma}

\begin{proof}
(M1)--(M3) hold for realized mosaics because $\mathcal{I}_w$ is a
strong abstract RCC5 frame (Section~\ref{sec:prelim}); they are
preserved under (SC1)--(SC4) closure operations.  (M4) holds by
Lemma~\ref{lem:model-safe}.  (M5)--(M6) hold by
Lemma~\ref{lem:status-partition} and Lemma~\ref{lem:open-cross}.
(M7) holds by the definition of the typed-$\EQ$ colouring
(Section~\ref{sec:typed-eq}).  (SC1)--(SC4) hold by construction
(we explicitly close).
\end{proof}


% ═══════════════════════════════════════════════════════════════
\section{Witness menus for $\DR/\PO$ existentials}
\label{sec:witness-menu}
% ═══════════════════════════════════════════════════════════════

For each rank-$d$ profile $\sigma \in S$ and each $\exists R.D \in
\lambda_Q(\sigma)$ with $R \in \{\DR, \PO\}$, the witness menu
records a target profile and relation.

\begin{construction}[Witness-menu extraction]\label{con:witness-menu}
For each $u \in T$ with $\st_d(u) = \sigma$ and each $\exists R.D
\in \lambda(u)$ with $R \in \{\DR,\PO\}$:
\begin{enumerate}
\item Locate the selected witness $b = w(\orig(u),\exists R.D) \in
  D_w(a_0)$ from the witness policy.
\item Find a tree occurrence $v$ with $\orig(v) = b$ (such $v$
  exists by Construction~\ref{con:split-tree}, possibly as a sibling
  occurrence under a common ancestor).
\item Record the witness-menu entry $(\exists R.D) \mapsto
  (\st_d(v), R)$ in $\Wit_d(\sigma)$.
\end{enumerate}
\end{construction}

\begin{lemma}[Witness-menu relation-safety]\label{lem:witness-safe}
Every entry $(\exists R.D) \mapsto (\sigma', R)$ in $\Wit_d(\sigma)$
satisfies: $\Need_R(\lambda_Q(\sigma)) \subseteq \lambda_Q(\sigma')$
and $\Need_{\inv(R)}(\lambda_Q(\sigma')) \subseteq
\lambda_Q(\sigma)$.
\end{lemma}

\begin{proof}
By Lemma~\ref{lem:model-safe} applied to the model edge
$(\orig(u),\orig(v))$ with relation $R$, and the fact that
$\lambda_Q(\sigma)$ depends only on $\sigma$ (the closure type is
the first profile component).
\end{proof}

\begin{remark}
$\PP$ and $\PPI$ existentials do \emph{not} appear in the witness
menu.  They are discharged by reachability over the parent function
of $Q$ (Section~\ref{sec:reachability}).  This separation is
deliberate: $\PP/\PPI$ witnesses are tracked structurally by the
tree backbone, while $\DR/\PO$ witnesses are tracked
combinatorially by the menu.
\end{remark}


% ═══════════════════════════════════════════════════════════════
\section{Typed-$\EQ$ congruence}
\label{sec:typed-eq}
% ═══════════════════════════════════════════════════════════════

The typed-$\EQ$ congruence is an explicit axiom on the quotient,
\emph{separate} from blocking equivalence and from the parent
function's self-loops.

\begin{definition}[Typed-$\EQ$ congruence]\label{def:typed-eq}
An equivalence relation $\eqcong$ on occurrences of $T$ (or on
profiles in $S$, projecting from $T$) is a \emph{typed-$\EQ$
congruence} when, for all $x, y$ with $x \eqcong y$:
\begin{enumerate}[label=(EQ\arabic*),leftmargin=2em]
\item $\lambda(x) = \lambda(y)$ (same closure type);
\item $c_{\EQ}(x) = c_{\EQ}(y)$ (same typed-$\EQ$ colour);
\item the vertical eventuality summaries (rank-$k$ profile
  components 7) at $x$ and $y$ agree;
\item for every other equivalence class $[z]$, the relation
  $\rho(\orig(x), \orig(z'))$ is the same for all $z' \in [z]$
  and matches $\rho(\orig(y), \orig(z'))$.
\end{enumerate}
\end{definition}

\begin{remark}[Blocking versus equality]\label{rem:blocking-vs-eq}
A profile loop $\mathsf{par}_Q(\sigma) = \sigma$
(Remark~\ref{rem:self-loops}) is \emph{not} a typed-$\EQ$
congruence on the unfolding: distinct unfolding occurrences of
$\sigma$ have the same profile but are \emph{not} $\EQ$-related.
The unfolding of $C\!\uparrow$ has $u_0 \PP u_1 \PP u_2 \PP \cdots$
with all $u_i$ distinct semantic occurrences sharing one profile;
their typed-$\EQ$ colours form a single class only because the
typed-$\EQ$ colouring assigns equal colours iff the elements are
genuine $\EQ$-mates in $\mathcal{I}_w$, not iff they have equal
profiles.
\end{remark}

\begin{lemma}[Extracted congruence]\label{lem:extracted-eq}
The relation on $T$ defined by $u \eqcong v$ iff $\orig(u) =
\orig(v)$ is a typed-$\EQ$ congruence in the sense of
Definition~\ref{def:typed-eq}.
\end{lemma}

\begin{proof}
(EQ1): same origin implies same type.  (EQ2): the typed-$\EQ$
colour is computed from the origin (a finite refinement of the
type, choosing one colour per $\EQ$-class in $\mathcal{I}_w$).
(EQ3): vertical summaries are computed from the origin's
ancestor/descendant relations in $\mathcal{I}_w$, which are
identical for $\EQ$-related origins.  (EQ4): if $\orig(u) =
\orig(v)$, then for any $z'$, $\rho(\orig(u),\orig(z')) =
\rho(\orig(v),\orig(z'))$ trivially.
\end{proof}


% ═══════════════════════════════════════════════════════════════
\section{Local coherence (C1)--(C5)}
\label{sec:coherence}
% ═══════════════════════════════════════════════════════════════

We verify the (C1)--(C5) coherence conditions
of~\cite[Definition~1.8]{SplitForest}, now restated for mosaics.

\begin{theorem}[Local coherence]\label{thm:coherence}
The descriptors $\Mos_k(s,t)$ extracted in
Construction~\ref{con:mosaic-extract} satisfy local coherence
(C1)--(C5).
\end{theorem}

\begin{proof}

\noindent\textbf{(C1) $\Out$ rigidity.}\;
By (M5) of Definition~\ref{def:mosaic}, every $\Out$-status
endpoint contributes only $\DR$-labelled cross-status edges.

\noindent\textbf{(C2) $\Front$--$\Front$ type-safety.}\;
Let $\mu \in \Mos_k(s,t)$ contain endpoints $v_1, v_2$ with
$(\Front,\Front)$ status.  By (M4),
$\mu(v_1,v_2) \in \{\DR,\PO\}$ is relation-safe for the closure
types $(\lambda(v_1), \lambda(v_2))$.

\noindent\textbf{(C3) Left downward support.}\;
Suppose $\mu \in \Mos_k(s,t)$ contains $v_1$ on the $s$-side, and
let $v_1'$ be a generated $\PPI$-child of $\orig(v_1)$
(equivalently, a child occurrence in $T_s$).  By (SC4) (one-step
containment extension), the descriptor $\Mos_k(s,t)$ contains a
mosaic $\mu'$ extending $\mu$ with $v_1'$ added.  The label
$\mu'(v_1', v_2)$ lies in $\comp(\PPI, \mu(v_1,v_2))$, hence the
required composition rule holds.

\noindent\textbf{(C4) Right downward support.}\;
Symmetric to (C3) on the $t$-side.

\noindent\textbf{(C5) Triple coherence.}\;
This is where v1's witness conflation occurred.  Let $v_1, v_2, v_3
\in V(\mu)$ be three vertices in a mosaic $\mu \in \Mos_k(s,t)$
with the relevant statuses.  Because $\mu$ is an \emph{atomic
mosaic} (a joint network on $V$, not a collection of independent
pair labels), it directly satisfies $\mu(v_1, v_3) \in
\comp(\mu(v_1, v_2), \mu(v_2, v_3))$ by (M3).

The crucial point: the three relations $\mu(v_1, v_2)$, $\mu(v_2,
v_3)$, $\mu(v_1, v_3)$ are realized \emph{by a single triple of
witnesses} $(\orig(v_1), \orig(v_2), \orig(v_3))$ in
$\mathcal{I}_w$, not by three independent pairs.  This is enforced
by Construction~\ref{con:mosaic-extract}: a realized mosaic comes
from a single vertex set $V$, with all pairs labelled by the
\emph{same} model.

Support closure (SC1)--(SC4) propagates this joint realizability:
restriction preserves it (a sub-triple of a real triple is a real
triple); typed-$\EQ$ replacement preserves it (an $\EQ$-mate is
realized by the same model element); one-step extension preserves
it (composition through the model).
\end{proof}

\begin{remark}[How v1's bug is repaired]
v1 wrote ``$M((\chi_1,\chi_2))$ contains $R_{12}$ from witnesses
$(x_1,x_2)$, $M((\chi_1,\chi_3))$ contains $R_{13}$ from witnesses
$(x_1', x_3)$, hence $\rho(\orig(x_2),\orig(x_3)) \in
\comp(\inv(R_{12}), R_{13})$ appears in $M((\chi_2,\chi_3))$.''  The
bug: $x_1$ and $x_1'$ may be distinct occurrences with the same
profile but different concrete witnesses; the composition does
\emph{not} apply across them.

The repair: mosaics record \emph{joint} networks on a single vertex
set.  Triple coherence is automatic because the three relations are
realized by one triple of model elements, by construction.
\end{remark}


% ═══════════════════════════════════════════════════════════════
\section{Reachability discharge for $\PP/\PPI$ eventualities}
\label{sec:reachability}
% ═══════════════════════════════════════════════════════════════

This section provides the finitary alternative to GPT-5.5's
cycle-realization condition~\cite[Definition~6.4]{GPT55Generated}.

\begin{definition}[Parent reachability in $Q$]\label{def:parent-reach}
For $\sigma, \sigma' \in S$, write $\sigma \reach_{\PP} \sigma'$ if
there exist $\sigma_0, \sigma_1, \ldots, \sigma_n$ in $S$ with
$\sigma_0 = \sigma$, $\sigma_n = \sigma'$, $n \ge 1$, and
$\mathsf{par}_Q(\sigma_i) = \sigma_{i+1}$ for $0 \le i < n$.
Self-loops are permitted: $\mathsf{par}_Q(\sigma) = \sigma$ implies
$\sigma \reach_{\PP} \sigma$.

Symmetrically, $\sigma \reach_{\PPI} \sigma'$ via the inverse
parent function (i.e., child function), allowing self-loops.
\end{definition}

\begin{lemma}[Reachability discharge]\label{lem:reach-discharge}
Let $Q$ be the quotient extracted in
Construction~\ref{con:split-tree} from a witness-generated model
$(\mathcal{I}_w, a_0)$.  For every $\sigma \in S$ and every
$\exists R.D \in \lambda_Q(\sigma)$ with $R \in \{\PP, \PPI\}$,
there exists $\sigma' \in S$ with $\sigma \reach_R \sigma'$ and
$D \in \lambda_Q(\sigma')$.
\end{lemma}

\begin{proof}
We treat $R = \PP$; $R = \PPI$ is symmetric.

Let $\sigma \in S$ with $\exists\PP.D \in \lambda_Q(\sigma)$.  Pick
any occurrence $u \in T$ with $\st_d(u) = \sigma$.  By
Construction~\ref{con:split-tree} step 2 (upward expansion), the
parent chain of $u$ contains an occurrence $u'$ with $D \in
\lambda(u')$ (the witness for $\exists\PP.D$).

Concretely: there is a finite parent chain $u =
u_0,u_1,\ldots,u_m$ in $T$ with $\mathsf{par}(u_i) = u_{i+1}$ and
$D \in \lambda(u_m)$.  Projecting to profiles via $\st_d$:
\[
\sigma = \st_d(u_0) \stackrel{\mathsf{par}_Q}{\to} \st_d(u_1)
\stackrel{\mathsf{par}_Q}{\to} \cdots
\stackrel{\mathsf{par}_Q}{\to} \st_d(u_m).
\]
This gives $\sigma \reach_{\PP} \sigma'$ with $\sigma' =
\st_d(u_m)$ and $D \in \lambda_Q(\sigma')$.
\end{proof}

\begin{remark}[This is finite-graph reachability, not B\"uchi]
\label{rem:not-buchi}
Lemma~\ref{lem:reach-discharge} reduces the discharge of
transitive $\PP/\PPI$ eventualities to ordinary reachability over
the finite directed graph $(S, \mathsf{par}_Q)$, decidable in
$O(|S| + |E_Q|)$ time per eventuality (e.g., breadth-first search).

GPT-5.5's Definition~6.4 phrases the condition in terms of stem-
and-cycle decompositions of branches, with realization required on
the cycle.  When the branch is summarized by a finite quotient $Q$,
``cycle in the unfolding'' becomes ``cycle in $\mathsf{par}_Q$''.
Realization on the cycle then becomes existence of a profile in the
strongly connected component containing $\sigma$ that has $D$ in
its type.

Restated as graph reachability: $\exists\PP.D \in \lambda_Q(\sigma)$
holds iff $D$ appears in some profile reachable from $\sigma$ via
$\mathsf{par}_Q$.  This is decidable in polynomial time without
B\"uchi acceptance, $\omega$-words, or any automata-theoretic
machinery.
\end{remark}

\begin{lemma}[Universal safety propagation]\label{lem:safety}
For every $\sigma \in S$ and every $\forall R.D \in
\lambda_Q(\sigma)$ with $R \in \{\PP, \PPI\}$, every $\sigma'$
with $\sigma \reach_R \sigma'$ has $D \in \lambda_Q(\sigma')$.
\end{lemma}

\begin{proof}
By Lemma~\ref{lem:model-safe} applied along the parent chain in
$\mathcal{I}_w$: every $\PP$-related pair satisfies the universal,
hence its target type contains $D$.  Profile equality preserves
type, so $\lambda_Q(\sigma') \ni D$.
\end{proof}


% ═══════════════════════════════════════════════════════════════
\section{Validity of the extracted quotient}
\label{sec:validity}
% ═══════════════════════════════════════════════════════════════

We verify that the quotient extracted in
Definition~\ref{def:quotient} is valid in the sense
of~\cite[Definition~1.10]{SplitForest} (or equivalently
of~\cite[Definition~8.2]{GPT55Generated}).

\begin{theorem}[Validity]\label{thm:validity}
The quotient $Q$ extracted from $(\mathcal{I}_w, a_0)$ in
Definition~\ref{def:quotient} satisfies all validity conditions:
\begin{enumerate}[label=(V\arabic*),leftmargin=2em]
\item every existential $\exists R.D \in \lambda_Q(\sigma)$ has a
  witness;
\item every universal $\forall R.D \in \lambda_Q(\sigma)$ is
  respected;
\item every domain in the quotient has non-empty $\Need_R$-filtered
  domain;
\item the extracted descriptors are support-closed and locally
  coherent;
\item $\eqcong$ is a typed-$\EQ$ congruence;
\item $C_0 \in \lambda_Q(\sigma_\ast)$.
\end{enumerate}
\end{theorem}

\begin{proof}

\noindent\textbf{(V1) Existential satisfaction.}\;
For $R \in \{\DR,\PO\}$: $\Wit_d(\sigma)$ provides
$(\sigma', R)$ with $D \in \lambda_Q(\sigma')$, by
Construction~\ref{con:witness-menu}.  Relation-safety holds by
Lemma~\ref{lem:witness-safe}.

For $R \in \{\PP, \PPI\}$: by Lemma~\ref{lem:reach-discharge},
some $\sigma'$ with $\sigma \reach_R \sigma'$ has $D \in
\lambda_Q(\sigma')$.

\noindent\textbf{(V2) Universal satisfaction.}\;
For $R \in \{\DR,\PO\}$: any sibling-mosaic edge from $\sigma$ to
$\sigma'$ labelled $R$ is realized by a model pair, which
satisfies the universal by Lemma~\ref{lem:model-safe}.

For $R \in \{\PP,\PPI\}$: Lemma~\ref{lem:safety} establishes
universal propagation along parent reachability.

\noindent\textbf{(V3) Non-empty filtered domains.}\;
For any pair-domain $\Dom(\sigma_1, \sigma_2)$ extracted from the
model:
\[
\Dom(\sigma_1,\sigma_2) = \{\rho(\orig(u),\orig(v)) : \st_d(u) =
\sigma_1, \st_d(v) = \sigma_2\}.
\]
Every $R \in \Dom(\sigma_1,\sigma_2)$ is realized by a model
pair, hence is relation-safe by Lemma~\ref{lem:model-safe}.
Therefore the $\Need_R$-filtered domain equals the unfiltered
domain, and is non-empty whenever $\Dom(\sigma_1,\sigma_2)$ is
non-empty.

\noindent\textbf{(V4) Coherent descriptors.}\;
By Theorem~\ref{thm:coherence}.

\noindent\textbf{(V5) Typed-$\EQ$ congruence.}\;
By Lemma~\ref{lem:extracted-eq}.

\noindent\textbf{(V6) Evaluation state.}\;
$u_\ast$ has $\orig(u_\ast) = a_0$ and $\lambda(u_\ast) =
\tp_{\mathcal{I}_w}(a_0) \ni C_0$.  Hence $C_0 \in
\lambda_Q(\sigma_\ast)$.
\end{proof}


% ═══════════════════════════════════════════════════════════════
\section{Main theorem and decidability}
\label{sec:main}
% ═══════════════════════════════════════════════════════════════

\begin{proof}[Proof of Theorem~\ref{thm:main}]
Given $\mathcal{I},d_\ast \models C_0$:
\begin{enumerate}
\item Choose a selected witness policy $w$
  (Definition~\ref{def:witness-policy}).
\item Pass to the witness-generated submodel $\mathcal{I}_w$
  (Lemma~\ref{lem:witness-generated}); $a_0 := d_\ast \in
  C_0^{\mathcal{I}_w}$.
\item Construct the generated split tree $T$
  (Construction~\ref{con:split-tree}, Lemma~\ref{lem:split-extraction}).
\item Define rank-$d$ profiles $\st_d$
  (Definition~\ref{def:profile-components}).
\item Form the finite quotient $Q$
  (Definition~\ref{def:quotient}).
\item Extract sibling descriptors as support-closed mosaics
  (Construction~\ref{con:mosaic-extract},
  Lemma~\ref{lem:descriptor-coherent}).
\item Extract witness menus
  (Construction~\ref{con:witness-menu},
  Lemma~\ref{lem:witness-safe}).
\item Verify validity (Theorem~\ref{thm:validity}).
\end{enumerate}
The resulting $Q$ is a valid finite generated split-forest quotient
for $C_0$ with $C_0 \in \lambda_Q(\sigma_\ast)$.
\end{proof}

\begin{proof}[Proof of Corollary~\ref{cor:decidability}]
$|\Sigma_d|$ is finite by Lemma~\ref{lem:finite-index}.  Validity
of a candidate quotient $Q$ is decidable: each (V1)--(V6) condition
is a finite check on finite data; reachability (V1) for $\PP/\PPI$
is decidable in $O(|S|^2)$ per eventuality; descriptor coherence
(V4) is decidable in time polynomial in $|\Mos|$.

If $C_0$ is satisfiable, Theorem~\ref{thm:main} produces a valid
quotient.  If a valid quotient exists, soundness
(\cite[Theorem~1.20]{SplitForest},
\cite[Theorem~9.1]{GPT55Generated}) produces a model.  Thus $C_0$
is satisfiable iff a valid finite quotient exists, and enumeration
terminates.
\end{proof}


% ═══════════════════════════════════════════════════════════════
\section{Discussion}
\label{sec:discussion}
% ═══════════════════════════════════════════════════════════════

\subsection{Why this is not B\"uchi acceptance}

The key observation that lets us avoid automata-theoretic
machinery: every transitive eventuality $\exists R.D$ at $\sigma
\in S$ for $R \in \{\PP, \PPI\}$ is discharged by ordinary
reachability over the finite directed graph $(S,
\mathsf{par}_Q)$.

In an unfolding of $Q$, an infinite $\PP$-branch corresponds to a
walk in $(S, \mathsf{par}_Q)$.  Once the walk enters a strongly
connected component (SCC), it cycles through that component
forever (because $|S|$ is finite).  Eventually-recurrent profiles
on the SCC are exactly the profiles reachable from any profile in
the SCC.

So: a $\PP$-eventuality $\exists\PP.D$ at $\sigma$ is realized in
the unfolding iff $D$ appears in some profile reachable from
$\sigma$.  This is decidable by:
\begin{enumerate}
\item BFS from $\sigma$ in $(S, \mathsf{par}_Q)$;
\item check whether any reachable profile $\sigma'$ has $D \in
  \lambda_Q(\sigma')$.
\end{enumerate}

GPT-5.5's stem-and-cycle phrasing
(``$u_0,\ldots,u_m,(v_0,\ldots,v_{k-1})^\omega$'',
\cite[Definition~6.4]{GPT55Generated}) is one specific way to
present this realization condition.  But the underlying check is
the same finite reachability question, and our formulation in
Lemma~\ref{lem:reach-discharge} is closer to standard description
logic blocking conditions.

\subsection{Comparison with v1}

\begin{center}
\renewcommand{\arraystretch}{1.25}
\begin{tabular}{p{0.32\linewidth}p{0.30\linewidth}p{0.30\linewidth}}
\toprule
Concern & v1 (April 2026) & v2 (this paper) \\
\midrule
Starting model & arbitrary $\PP$-DAG with assumed $\le_{\PP}$-min &
witness-generated submodel \\
Cover edges & Hasse covers of ambient order & generated cover
edges in $D_w(a_0)$ \\
Ambient root & required (added freshly if needed) & not required \\
Self-loops on parent & not permitted & permitted (handles
$C\!\uparrow$) \\
Sibling descriptor & raw pair-tables & support-closed mosaics \\
Triple coherence (C5) & witness conflation across occurrences &
joint-realization within a mosaic \\
Typed-$\EQ$ congruence & implicit via $\orig$ map & explicit axiom
\\
$\PP/\PPI$ eventualities & ``ancestor in tree'' (vague) & finite
graph reachability \\
\bottomrule
\end{tabular}
\end{center}

\subsection{Comparison with GPT-5.5's automata version}

GPT-5.5's coherent generated split-forest
proof~\cite{GPT55Generated} arrives at the same conclusion via the
same structural ideas (witness-generated submodels, generated
covers, support-closed mosaics, typed-$\EQ$ congruence) but
formalizes the eventuality discharge as a B\"uchi-style
cycle-realization condition.

Our v2 differs only in presentation: where GPT-5.5 writes
$u_0\cdots u_m(v_0\cdots v_{k-1})^\omega$ and requires every
eventuality on the cycle to be met on the cycle, we write
``every $\exists\PP.D$ at $\sigma$ has a reachable witness
profile,'' which is finite-graph reachability.  These are
equivalent: in the finite quotient, ``met somewhere on the
$\omega$-trace'' equals ``reachable in $\mathsf{par}_Q$.''

The benefit of the reachability formulation: it stays inside the
standard description-logic toolbox (similar to fixpoint blocking
conditions in $\mu$-calculus tableaux), and avoids invoking
$\omega$-language theory.

\subsection{Honest assessment}

We flag the following points for scrutiny:

\paragraph{Witness policy is non-canonical.}
Different selected witness policies produce different generated
submodels and different split trees.  The extraction works for
\emph{any} fixed policy.  We do not claim a canonical choice.

\paragraph{Profile bound is not elementary.}
$|\Sigma_d|$ is bounded by a tower of exponentials in $|\cl(C_0)|$.
Decidability follows, but tight complexity remains open.  The
PO-coherent fragment achieves a 2-EXPTIME quotient
bound~\cite{TwoTier}.

\paragraph{Mosaic vertex bound.}
Construction~\ref{con:mosaic-extract} restricts to mosaics on
$\le |\cl(C_0)| + 4$ vertices.  This is sufficient for
support-closure (SC1)--(SC4), but a fully rigorous treatment
would derive the bound from a Ramsey-type argument on closure
size.  GPT-5.5 sketches such an argument
in~\cite[Lemma~8.4]{GPT55Generated}.

\paragraph{Reachability vs strongly-connected components.}
Lemma~\ref{lem:reach-discharge} establishes existence of a
reachable witness profile for each eventuality.  The full
correctness argument requires that \emph{simultaneous} discharge
of all eventualities on a branch is realized by a single coherent
infinite walk.  This follows from finiteness of $S$ and the
RCC5 patchwork property: the SCC of $\sigma$ in $(S,\mathsf{par}_Q)$
contains all profiles eventually visited from $\sigma$, and any
ordering of the eventualities can be interleaved on the SCC walk.

None of these points appears to harbour a genuine gap.  The
extraction rests on three clean principles:
\begin{enumerate}
\item witness-generated submodels are sparse
  (Lemma~\ref{lem:witness-generated});
\item model edges are relation-safe
  (Lemma~\ref{lem:model-safe});
\item finite-graph reachability discharges transitive eventualities
  (Lemma~\ref{lem:reach-discharge}).
\end{enumerate}


% ═══════════════════════════════════════════════════════════════
\section{Conclusion}
\label{sec:conclusion}
% ═══════════════════════════════════════════════════════════════

We have written out the completeness extraction for $\ALCIRCC{5}$
in a repaired form that fixes the two structural defects of the
v1 draft: the unsound Hasse construction and the witness-conflated
triple-coherence proof.  The repair uses witness-generated
submodels, generated cover edges, support-closed mosaics, an
explicit typed-$\EQ$ congruence, and finite-graph reachability
discharge.

Crucially, the repair stays inside the original split-forest
framework: vertical $\PP/\PPI$ branches, sibling-interface
descriptors, and finite quotient enumeration are all preserved.
The cycle-realization condition that motivates GPT-5.5's
automata-flavoured proof reduces here to ordinary reachability
over the finite quotient graph, sidestepping any need for B\"uchi
acceptance or $\omega$-language theory.

Combined with the soundness chain
of~\cite{SplitForest,GPT55Generated}, this yields:

\begin{center}
\emph{Concept satisfiability in $\ALCIRCC{5}$ under strong $\EQ$
semantics is decidable.}
\end{center}

The exact complexity is open.  The quotient search is decidable but
non-elementary; an $\textsc{ExpTime}$ bound (matching the known
$\mathcal{ALCI}$ lower bound) would require a tighter analysis
of the rank-$d$ profile space.

\medskip
\noindent\textbf{Computational evidence.}\;
An independent implementation~\cite{ClaudeImpl} cross-validates
the cover-tree tableau (the operational counterpart of the
split-forest quotient search) against a quasimodel-based reasoner
on 911 test concepts with zero mismatches.  A decomposition test
confirms that 775/775 satisfiable concepts (100\%) have finite
cover-tree models across 12.2\,M enumerated models, with zero
genuine counterexamples~\cite{StressTestCoverTree,
DecompositionTest}.


\begin{thebibliography}{9}

\bibitem{Wessel2001}
M.~Wessel.
\newblock Obstacles on the way to qualitative spatial reasoning with
  description logics.
\newblock In \emph{Proc.\ DL Workshop}, 2001.

\bibitem{Wessel2002}
M.~Wessel.
\newblock On spatial reasoning with description logics---position paper.
\newblock In \emph{Proc.\ DL Workshop}, 2002.

\bibitem{Wessel2003}
M.~Wessel.
\newblock Qualitative spatial reasoning with the
  {$\mathcal{ALCI}_\mathit{RCC}$} family---{F}irst results and
  unanswered questions.
\newblock Technical Report FBI-HH-M-324/03, University of Hamburg, 2002/2003.

\bibitem{SplitForest}
GPT-5.4 (OpenAI).
\newblock Finite sibling-interface descriptors and a proposed completion
  for $\mathcal{ALCI}_{RCC5}$.
\newblock Proof-style note, April 2026.
\newblock Prompted by Michael Wessel.

\bibitem{GPT55Generated}
GPT-5.5 (OpenAI).
\newblock A coherent generated split-forest decidability proof for
  $\mathcal{ALCI}_{RCC5}$.
\newblock Proof manuscript, May 2026.
\newblock Prompted by Michael Wessel.

\bibitem{GPT55Survival}
GPT-5.5 (OpenAI).
\newblock What survives of the split-forest idea.
\newblock Companion note to the repaired decidability proof,
  May 2026.
\newblock Prompted by Michael Wessel.

\bibitem{CompletenessV1}
Claude (Anthropic).
\newblock Completeness extraction for $\mathcal{ALCI}_{RCC5}$ (v1).
\newblock Proof-style note, April 2026.
\newblock Prompted by Michael Wessel.
\newblock Superseded by the present paper.

\bibitem{TwoTier}
Claude (Anthropic).
\newblock Two-tier quotient decidability for the PO-coherent fragment
  of $\mathcal{ALCI}_{RCC5}$.
\newblock Proof-style note, March 2026.
\newblock Prompted by Michael Wessel.

\bibitem{ClaudeImpl}
Claude (Anthropic).
\newblock Cover-tree tableau implementation for $\mathcal{ALCI}_{RCC5}$.
\newblock Implementation paper, April 2026.
\newblock Prompted by Michael Wessel.

\bibitem{StressTestCoverTree}
Claude (Anthropic).
\newblock Stress test cross-validation suite (911 concepts, 0
  mismatches).
\newblock Python implementation, April 2026.

\bibitem{DecompositionTest}
Claude (Anthropic).
\newblock Cover-tree decomposition test (775/775 = 100\%, 12.2\,M
  models).
\newblock Python implementation, April 2026.

\end{thebibliography}

\end{document}
